\documentclass[10pt]{mypackage}

\usepackage[uprightgreeks,noDcommand]{kpfonts}
\renewcommand{\coloneq}{\coloneqq}
\renewcommand{\eqcolon}{\eqqcolon}
%\usepackage{dsfont}
%\renewcommand{\mathbb}{\mathds}
%\usepackage{homework}
\usepackage{notes}

%\usepackage[ backend=bibtex, style = alphabetic, sorting=ynt ]{biblatex}
%\addbibresource{  }

\usepackage{parskip}
\frenchspacing

\fancyhf{}
\fancyhead[R]{Avinash Iyer}
\fancyhead[L]{My Personal Style Guide}
\fancyfoot[C]{\thepage}

\setcounter{secnumdepth}{0}

\begin{document}
\RaggedRight
This is my personal style guide/notation guide as of 2026. I am writing this primarily for my own benefit so I know what notation I will use as I comb through sources.
\begin{itemize}
  \item Sets are denoted by standard Roman characters, such as $A,B,C$, or capitalized Greek characters such as $\Omega,\Lambda,\Gamma$.
  \item Vector spaces are denoted by \verb|\mathscr|, such as $ \mathscr{H} $ or $\mathscr{X}$.
  \item Nets or sequences are denoted by $\left( x_i \right)_{i\in I}$ or $\left( x_n \right)_{n\in \N}$.
  \item Indexed sets are written $\set{x_i : i\in I}$ or $\set{x_n : n\in \N}$.
  \item Families of subsets (such as topologies or otherwise) are denoted by \verb|\mathcal|, such as $\mathcal{T}$ or $\mathcal{U}$.
  \item The system of open neighborhoods of a point $y\in X$ is denoted $ \mathcal{O}_{X}(y) $. The system of all neighborhoods of a point $y\in X$ is denoted $\mathcal{N}_{X}(y)$.
  \item All vector spaces are complex unless stated otherwise.
  \item If $\left( X,d \right)$ is a metric space, then 
    \begin{align*}
      U\left( x,\ve \right) &= \set{y\in X : d\left( x,y \right) < \ve}\\
      B\left( x,\ve \right) &= \set{y\in X : d\left( x,y \right) \leq \ve}\\
      S\left( x,\ve \right) &= \set{y\in X : d\left( x,y \right) = \ve}.
    \end{align*}
  \item The Roman lettering $L^p$ spaces are used to denote almost-everywhere equivalence classes of $p$-integrable functions, while the \verb|\mathcal| lettering $ \mathcal{L}^p $ spaces are used to denote the space of all the $p$-integrable functions without considering almost-everywhere equivalence.
  \item If $ \mathscr{X} $ is a normed vector space, then $ \operatorname{ball}\left( \mathscr{X} \right) $ is used to denote the \textit{closed} unit ball of $ \mathscr{X} $ about the origin.
  \item If $ \mathscr{X} $ is a normed vector space, then $ \mathcal{B}\left( \mathscr{X} \right) $ denotes the space of bounded linear operators from $\mathscr{X}$ to itself.
  \item If $\mathscr{H}$ is a Hilbert space, then $ \mathcal{U}\left( \mathscr{H} \right) $ is the space of unitary operators on $ \mathscr{H} $.
  \item Operators in $ \mathcal{B}\left( \mathscr{H} \right) $ for some Hilbert space $ \mathscr{H} $ are denoted by standard Roman characters, such as $T,S,U$.
  \item If $ \mathscr{K}\subseteq \mathscr{H} $ is a closed subspace, then $\proj\left( \mathscr{K} \right)$ is used to denote the projection onto the subspace $ \mathscr{K} $.
  \item If $p$ and $q$ are two projections, then $p\vee q$ is their supremum and $p\wedge q$ is their infimum. If $p$ and $q$ are orthogonal to each other, then we write $p\vee q = p + q$.
  \item $C^{\ast}$-Algebras or von Neumann algebras will be written with the \verb|\mathcal| lettering, such as $ \mathcal{A} $ or $\mathcal{M}$.
  \item If $A\subseteq \mathcal{B}\left( \mathscr{H} \right)$ is any subset of operators, then $W^{\ast}\left( A \right)$ is used to denote the von Neumann algebra generated by the set $A$.
  \item If $ \mathcal{M} $ is a von Neumann algebra, then $\Proj\left( \mathcal{M} \right)$ is used for the projection lattice of $\mathcal{M}$.
\end{itemize}
\end{document}
